\documentclass{beamer}

\mode<presentation>
{
  \usetheme{Singapore}      % or try Darmstadt, Madrid, Warsaw, ...
  \usecolortheme{crane} % or try albatross, beaver, crane, ...
  \usefonttheme{serif}  % or try serif, structurebold, ...
  \setbeamertemplate{navigation symbols}{}
  \setbeamertemplate{caption}[numbered]
} 

\usepackage[autostyle=true]{csquotes}
\usepackage{xeCJK}
\usepackage[T1]{fontenc}
\usepackage{xcolor}
\usepackage{setspace}
\setmainfont[AutoFakeSlant=0.2]{Charis SIL}
\setCJKmainfont[AutoFakeSlant=0.2]{Noto Serif CJK HK}
\setCJKmonofont[AutoFakeSlant=0.2]{Noto Serif CJK HK}
\usepackage{polyglossia}
\newfontfamily\thaifont[Script=Thai]{Noto Serif Thai}
\setmainlanguage{english}
\setotherlanguages{thai}
\usepackage{csquotes}
\usepackage{array}
\usepackage{makecell}
\usepackage{changepage}
\usepackage{textcomp}
\usepackage{multirow}

\title[Pearl River Sinitic]{Yuè \& the other Sinitic Languages of the Pearl River Region --- Unit 5.2}
\subtitle{13ᵗʰ EACL Summer School, Universität Stuttgart}
\author{Hilário de Sousa 蘇沙}
\institute{EHESS–CRLAO \\ {\tiny membre associé}}
\date{Thursday 31 July 2026 \\ 丙午六月十八 \\ {\footnotesize \texttt{hilario@bambooradical.com}}}

\begin{document}

\begin{frame}
  \titlepage
\end{frame}

\section{SW Mandarin etc.}

\begin{frame}[fragile]{SW Mandarin etc.}

\centering
\includegraphics[width=0.85\linewidth]{langatlaschinaguangxisinitic}\\
{\tiny Language Atlas of China 1\textsuperscript{st} ed.}\\
{\footnotesize Mandarin speakers mostly arrived in late 14\textsuperscript{th}/ early 15\textsuperscript{th} C (元末 / 明初)}\\

\end{frame}

\begin{frame}[fragile]{SW Mandarin etc.}

桂柳話 (話 of 桂林 and 柳州)\\
Guìlín accent: kwɛj˨˦ (nin˨˩) njəw˥˦ (tsəw˧) xwa˨˦\\
\quad \textit{\footnotesize more influences from Xiāng}\\
\quad {\footnotesize Guìlín-type Mandarin also spoken in S Húnán / border N Guangdong}\\
\quad {\footnotesize Guìlín: ˧ ˨˩ ˥˦ ˨˦}\\
\medskip
Liǔzhōu accent: kwɐj˨˦ (lən˧˩) ljɐw˥˦ (tsɐw˦) hwa˨˦\\
\quad \textit{\footnotesize more influences from Cantonese and N Zhuàng}\\
\quad {\footnotesize Liǔzhōu: ˦ ˧˩ ˥˦ ˨˦ (˥ʔ)}\\
{\footnotesize Guìliǔ-type Mandarin also spoken in S Guìzhōu}\\
{\footnotesize Usually no retroflex initials, often has [y]}\\
\bigskip
雲南話\\
Kūnmíng accent: jĩ˧˩ nã˧˩ xwa˨˩˨\\
\quad {\footnotesize Kūnmíng: ˦ ˧˩ ˥ ˨˩˨}\\
{\footnotesize Often has retroflex initials (except Far-S), usually no [y]}

\end{frame}

\begin{frame}[fragile]{SW Mandarin etc.}

Liǔzhōu Mandarin 柳州官話\\
\bigskip
p pʰ m f, t tʰ n l, ts tsʰ s (tɕ tɕʰ ɕ), k kʰ ŋ h\\
{\footnotesize (ts tsʰ s > tɕ tɕʰ ɕ / \_{}\_{} i u j w)}\\
\bigskip
\includegraphics[width=0.9\linewidth]{liuzhourimes}

\end{frame}

\begin{frame}[fragile]{SW Mandarin etc.}

\begin{tabular}{l||c|c|c||c}
    \hline
    Standard Mand. & *平 *A & *上 *B & *去 *C & *入 *D\\\hline
    陰 {\tiny *voiceless 清 initial} & 1 ˥ & 3 ˨˩˥ & 4 ˥˩ & \textit{\footnotesize random}\\ \cline{1-2} \cline{5-5} 
    陽 {\tiny *voiced sonorant 次濁 initial} & 2 ˨˥ & & & 4 ˥˩\\ \cline{3-3} \cline{5-5}
    陽 {\tiny *voiced obstruent 全濁 initial} & & \multicolumn{1}{|l}{} & & 2 ˨˥\\ \hline
\end{tabular}

\bigskip
\begin{tabular}{l||c|c|c||c}
    \hline
    柳州官話 & *平 *A & *上 *B & *去 *C & *入 *D\\\hline
    陰 {\tiny *voiceless 清 initial} & 1 ˦ & 3 ˥˦ & 4 ˨˦ & 2 ˧˩\\ \cline{1-2} 
    陽 {\tiny *voiced sonorant 次濁 initial} & 2 ˧˩ & & & \\ \cline{3-3}
    陽 {\tiny *voiced obstruent 全濁 initial} & & \multicolumn{1}{|l}{} & & \\ \hline
\end{tabular}
{\footnotesize (Plus ˥ʔ borrowed from Cantonese)}\\
tone D > tone A2 very common amongst SW Mandarin dialects

\bigskip
一 i˧˩ 六 lu˧˩ 七 tɕʰi˧˩ 八 pa˧˩ 十 sz˧˩\\
豛 toʔ˥ < Canto. duk1 \enquote*{poke}\\
嘜 maʔ˥ < Canto. mak1 \enquote*{brand} (< Eng. \enquote{mark})\\

\end{frame}

\begin{frame}[fragile]{SW Mandarin etc.}

{\footnotesize (Based on Tang 2012)}\\
\medskip
\centering
    {\small
    \begin{tabular}{l l}
    %\begin{tabular}{p{5cm} p{5cm}}
        \textbf{柳州官話} & \textbf{桂林官話}\\ \hline
        泥來不混 n ≠ l & 泥來相混 n \~{} l\\
        泥 ni˧˩ ≠ 梨 li˧˩ & 泥 = 梨 ni˨˩\\ \hline
        分尖團 MC tsj ≠ kj & 不分尖團 MC tsj = kj\\
        精 tsən˦ ≠ 京 kjən˦ & 精 = 京 tɕin˧\\ \hline
        公雞 / 雞公 & 公雞 \\ 
        koŋ˦ ki˦ / ki˦ koŋ˦ & koŋ˧ tɕi˧\\ \hline
        你們坐先 & 先和他講桂林話\\ 
        ni˥˦ mɐn˧˩ tso˨˦ sẽ˦ & ɕjɛ̃˧ xo˨˩ tʰa˧ tɕjaŋ˥˦ kwɛj˨˦ nin˨˩ xwa˨˦\\ \hline
    \end{tabular}}

\vfill \setstretch{0.5} {\tiny Táng, Qīyuán 唐七元. 2012. 试谈桂林官话和柳州官话之间的方言差异. Journal of Guilin University of Aerospace Technology 桂林航天工业学院学报 2012(4). 453–456.}\par

\end{frame}

\begin{frame}[fragile]{SW Mandarin etc.}

Example of other (Ming-dynasty) Mandarin enclaves, e.g.\\
邕州官話 Yōngzhōu Guānhuà (Old Nánníng Mandarin)\\
\medskip
p pʰ m f, t tʰ n l, ts tsʰ s (tɕ tɕʰ ɕ) ɬ, k kʰ ŋ h\\

\includegraphics[width=0.9\linewidth]{nanningmandrimes}\\
˧˥ ˧˩ ˥˦ ˩˧ / ˥ ˨\\
\end{frame}

\begin{frame}[fragile]{SW Mandarin etc.}

邕州官話 Yōngzhōu Guānhuà (Old Nánníng Mandarin)\\
˧˥ ˧˩ ˥˦ ˩˧ / ˥ ˨\\
\medskip
(SW) Mandarin traits: e.g.\\
*-m > *-n, *-ŋ \textsuperscript{sometimes} > -n\\
金 = 斤 = 京 kin˧˥ \quad 針 = 珍 = 貞 tɕin˧˥\\ 慘 = 產 = 橙 tsʰan˥˦\\
\smallskip
*-p > *-t > *-k > -0 (or glide)\\
濕 = 實 = 石 ɕi˧˩\\
\medskip
Quite many loans from e.g. Nan. Canto.\\
鉗 kʰjem˧˩ (< NC kʰim˨˩), 粒 nəp˥ (< NC nɐp˥)\\ 
忽 fɐt˥ (< NC fɐt˥), 突 tət˨ (< NC tɐt˨)\\
鑿 tsɔk˨ (< NC tsɔk˨)

\end{frame}

\begin{frame}[fragile]{SW Mandarin etc.}

「軍話」\enquote{Jūnhuà}:\\
Many enclaves of descendants of Ming-dynasty soldiers, speaking some type of \enquote{older} Mandarin.\\
(The soldiers may or may not be native speakers of Mandarin; some kind of koine Mandarin developed within each garrison.)\\

\end{frame}

\begin{frame}[fragile]{SW Mandarin etc.}

e.g. 舊時正話 of 電白 (Chen 2006)\\
Ming-dynasty soldiers primarily from Zhejiang or Fujian (i.e. none or very few of them were native speakers of Mandarin)\\
電白華樓 A1 ˧ A2 ˨˩˧ B ˧˩ C ˥ / DL ˨˩˧ DS ˥\\
\medskip
\includegraphics[width=1\linewidth]{dianbaipronouns}\\
\bigskip
The 電白舊時正話 community does not have the endonym/exonym of 軍話;\\
For Qiū (2005: 136), only Mandarin varieties which have the endonym or exonym of 軍話 should be considered 軍話.\\

\vfill \setstretch{0.5} {\tiny Chén, Yúnlóng 陈云龙. 2006. 旧时正话研究. Beijing: China Social Science Press 中国社会科学出版社.\\
Qiū, Xuéqiáng 丘学强. 2005. 军话研究. Beijing: China Social Science Press 中国社会科学出版社.}\par

\end{frame}

\begin{frame}[fragile]{SW Mandarin etc.}

Notice: not all communities of descendants of such soldiers speak a Mandarin variety\\
e.g. 大鵬話 of Shenzhen/HK is mixed Hakka/Yue\\

\end{frame}

\begin{frame}[fragile]{SW Mandarin etc.}

Manchu quarters 滿城 (e.g. Canton 1860:)\\
\includegraphics[width=1\linewidth]{canton1860}

\end{frame}

\begin{frame}[fragile]{SW Mandarin etc.}
廣州：旗下話 (NE Mandarin) \& 滿洲話 (Beijing Mandarin)\\
both with Cantonese influences; moribund\\
\bigskip
滿洲話 (Yang 2023) e.g.:\\
ku˥ tu˥ \enquote*{tomb} < Manchu \textit{huru} \enquote*{high place} / \textit{kuru} \enquote*{elevation}\\
擦黑{\scriptsize 儿} \enquote{dusk, nightfall} (e.g. NE Mandarin)\\
\medskip
Tons of Cantonese calques, e.g.:\\
幾多 jǐduō \enquote*{how/so many} < Canto. gei2 do1\\
唔該晒 bùgāi saai3 \enquote{thanks} < Canto. m4 goi1 saai3\\
隔籬 gélí \enquote*{neighbouring} < Canto. gaak3 lei4\\

\vfill \setstretch{0.5} {\tiny Yang, Gareth Junjie 楊鈞傑. 2023. The Moribund “Manchu Mandarin” in Guangzhou: A Product of Cantonese-Mandarin Contact. Paper presented at The 27th International Conference on Yue Dialects, The Ohio State University.}\par

\end{frame}

\begin{frame}[fragile]{SW Mandarin etc.}

廣州滿洲話 (Yang 2023) e.g.:\\

\bigskip
包饽饽要买啲灰面和豆豉酱\\
bāo bōbo yào mǎi di1 huīmiàn hé dòuchǐ-jiàng\\
\enquote{To make dumplings, [you] need to buy some flour and \textit{douchi} sauce}\\

\medskip
啲 di1: Canto. mass classifier

\bigskip
你在哪啲找吃？\\
nǐ zài nǎdi1 zhǎo chī\\
\enquote{Where do you make your living?}\\

\medskip
找吃 < Canto. 搵食 wan2 sik6 (seek eat) \enquote*{make a living}\\
哪啲：\\哪裏 nali > nadi(?); influence of Canto. mass \textsc{clf} 啲 di1(?)\\ 

\end{frame}

\begin{frame}[fragile]{SW Mandarin etc.}

粵劇官話 （戲棚官話／中州音）\\
Cantonese opera were all sung in Mandarin in earlier days\\
Now primarily in Cantonese, but some Mandarin is still used\\
\medskip 
Cantonese Opera Mandarin is like earlier Guilin Mandarin + Cantonese influence, e.g. \url{https://www.canoperadict.com/}\\
\medskip
Cantonese Opera Mandarin words in Cantonese: e.g.\\
可怒也！ ko1 naau1 je1!\\
合晒合尺 hap6 saai3 ho4 ce1\\

\bigskip
{\footnotesize
\begin{tabular}{c c c c c c c c c c c}
    ... & 合 & 士 & 乙 & 上 & 尺 & 工 & 反 & 六 & 五 & ... \\
    ... & ho4 & si6 & ji6 & saang3 & ce1 & gung1 & faan1 & liu1 & wu1 & ...\\
    ... & G & A & B & C & D & E & F\# & G & A & ...\\
    & \textcolor{black!40}{hap6} & \textcolor{black!40}{si6} & \textcolor{black!40}{jyut6} & \textcolor{black!40}{soeng6} & \textcolor{black!40}{cek3} & \textcolor{black!40}{gung1} & \textcolor{black!40}{faan2} & \textcolor{black!40}{luk6} & \textcolor{black!40}{ng5} & \\
\end{tabular}}

\end{frame}


\end{document}
